% Options for packages loaded elsewhere
\PassOptionsToPackage{unicode}{hyperref}
\PassOptionsToPackage{hyphens}{url}
%
\documentclass[
]{article}
\usepackage{amsmath,amssymb}
\usepackage{iftex}
\ifPDFTeX
  \usepackage[T1]{fontenc}
  \usepackage[utf8]{inputenc}
  \usepackage{textcomp} % provide euro and other symbols
\else % if luatex or xetex
  \usepackage{unicode-math} % this also loads fontspec
  \defaultfontfeatures{Scale=MatchLowercase}
  \defaultfontfeatures[\rmfamily]{Ligatures=TeX,Scale=1}
\fi
\usepackage{lmodern}
\ifPDFTeX\else
  % xetex/luatex font selection
\fi
% Use upquote if available, for straight quotes in verbatim environments
\IfFileExists{upquote.sty}{\usepackage{upquote}}{}
\IfFileExists{microtype.sty}{% use microtype if available
  \usepackage[]{microtype}
  \UseMicrotypeSet[protrusion]{basicmath} % disable protrusion for tt fonts
}{}
\makeatletter
\@ifundefined{KOMAClassName}{% if non-KOMA class
  \IfFileExists{parskip.sty}{%
    \usepackage{parskip}
  }{% else
    \setlength{\parindent}{0pt}
    \setlength{\parskip}{6pt plus 2pt minus 1pt}}
}{% if KOMA class
  \KOMAoptions{parskip=half}}
\makeatother
\usepackage{xcolor}
\usepackage[margin=1in]{geometry}
\usepackage{color}
\usepackage{fancyvrb}
\newcommand{\VerbBar}{|}
\newcommand{\VERB}{\Verb[commandchars=\\\{\}]}
\DefineVerbatimEnvironment{Highlighting}{Verbatim}{commandchars=\\\{\}}
% Add ',fontsize=\small' for more characters per line
\usepackage{framed}
\definecolor{shadecolor}{RGB}{248,248,248}
\newenvironment{Shaded}{\begin{snugshade}}{\end{snugshade}}
\newcommand{\AlertTok}[1]{\textcolor[rgb]{0.94,0.16,0.16}{#1}}
\newcommand{\AnnotationTok}[1]{\textcolor[rgb]{0.56,0.35,0.01}{\textbf{\textit{#1}}}}
\newcommand{\AttributeTok}[1]{\textcolor[rgb]{0.13,0.29,0.53}{#1}}
\newcommand{\BaseNTok}[1]{\textcolor[rgb]{0.00,0.00,0.81}{#1}}
\newcommand{\BuiltInTok}[1]{#1}
\newcommand{\CharTok}[1]{\textcolor[rgb]{0.31,0.60,0.02}{#1}}
\newcommand{\CommentTok}[1]{\textcolor[rgb]{0.56,0.35,0.01}{\textit{#1}}}
\newcommand{\CommentVarTok}[1]{\textcolor[rgb]{0.56,0.35,0.01}{\textbf{\textit{#1}}}}
\newcommand{\ConstantTok}[1]{\textcolor[rgb]{0.56,0.35,0.01}{#1}}
\newcommand{\ControlFlowTok}[1]{\textcolor[rgb]{0.13,0.29,0.53}{\textbf{#1}}}
\newcommand{\DataTypeTok}[1]{\textcolor[rgb]{0.13,0.29,0.53}{#1}}
\newcommand{\DecValTok}[1]{\textcolor[rgb]{0.00,0.00,0.81}{#1}}
\newcommand{\DocumentationTok}[1]{\textcolor[rgb]{0.56,0.35,0.01}{\textbf{\textit{#1}}}}
\newcommand{\ErrorTok}[1]{\textcolor[rgb]{0.64,0.00,0.00}{\textbf{#1}}}
\newcommand{\ExtensionTok}[1]{#1}
\newcommand{\FloatTok}[1]{\textcolor[rgb]{0.00,0.00,0.81}{#1}}
\newcommand{\FunctionTok}[1]{\textcolor[rgb]{0.13,0.29,0.53}{\textbf{#1}}}
\newcommand{\ImportTok}[1]{#1}
\newcommand{\InformationTok}[1]{\textcolor[rgb]{0.56,0.35,0.01}{\textbf{\textit{#1}}}}
\newcommand{\KeywordTok}[1]{\textcolor[rgb]{0.13,0.29,0.53}{\textbf{#1}}}
\newcommand{\NormalTok}[1]{#1}
\newcommand{\OperatorTok}[1]{\textcolor[rgb]{0.81,0.36,0.00}{\textbf{#1}}}
\newcommand{\OtherTok}[1]{\textcolor[rgb]{0.56,0.35,0.01}{#1}}
\newcommand{\PreprocessorTok}[1]{\textcolor[rgb]{0.56,0.35,0.01}{\textit{#1}}}
\newcommand{\RegionMarkerTok}[1]{#1}
\newcommand{\SpecialCharTok}[1]{\textcolor[rgb]{0.81,0.36,0.00}{\textbf{#1}}}
\newcommand{\SpecialStringTok}[1]{\textcolor[rgb]{0.31,0.60,0.02}{#1}}
\newcommand{\StringTok}[1]{\textcolor[rgb]{0.31,0.60,0.02}{#1}}
\newcommand{\VariableTok}[1]{\textcolor[rgb]{0.00,0.00,0.00}{#1}}
\newcommand{\VerbatimStringTok}[1]{\textcolor[rgb]{0.31,0.60,0.02}{#1}}
\newcommand{\WarningTok}[1]{\textcolor[rgb]{0.56,0.35,0.01}{\textbf{\textit{#1}}}}
\usepackage{graphicx}
\makeatletter
\newsavebox\pandoc@box
\newcommand*\pandocbounded[1]{% scales image to fit in text height/width
  \sbox\pandoc@box{#1}%
  \Gscale@div\@tempa{\textheight}{\dimexpr\ht\pandoc@box+\dp\pandoc@box\relax}%
  \Gscale@div\@tempb{\linewidth}{\wd\pandoc@box}%
  \ifdim\@tempb\p@<\@tempa\p@\let\@tempa\@tempb\fi% select the smaller of both
  \ifdim\@tempa\p@<\p@\scalebox{\@tempa}{\usebox\pandoc@box}%
  \else\usebox{\pandoc@box}%
  \fi%
}
% Set default figure placement to htbp
\def\fps@figure{htbp}
\makeatother
\setlength{\emergencystretch}{3em} % prevent overfull lines
\providecommand{\tightlist}{%
  \setlength{\itemsep}{0pt}\setlength{\parskip}{0pt}}
\setcounter{secnumdepth}{-\maxdimen} % remove section numbering
\usepackage{bookmark}
\IfFileExists{xurl.sty}{\usepackage{xurl}}{} % add URL line breaks if available
\urlstyle{same}
\hypersetup{
  pdftitle={ActividadFinal1},
  pdfauthor={OscarMacielC},
  hidelinks,
  pdfcreator={LaTeX via pandoc}}

\title{ActividadFinal1}
\author{OscarMacielC}
\date{2025-08-11}

\begin{document}
\maketitle

\pandocbounded{\includegraphics[keepaspectratio]{https://www.oscarmacielc.com/_next/image?url=\%2F_next\%2Fstatic\%2Fmedia\%2FOscarLogo.cdfce2e5.png&w=64&q=75}}
\href{https://github.com/OscarMacielC/RHerramientasProductividad}{El
github para la clase}
\pandocbounded{\includegraphics[keepaspectratio]{https://www.oscarmacielc.com/_next/image?url=\%2F_next\%2Fstatic\%2Fmedia\%2FOscarLogo.cdfce2e5.png&w=64&q=75}}

\begin{center}\rule{0.5\linewidth}{0.5pt}\end{center}

\section{Actividad 1}\label{actividad-1}

\subsection{1 Explora el entorno de
RStudio:}\label{explora-el-entorno-de-rstudio}

\textbf{¿Qué ves en los paneles?} varias ventanas para hacer cosas,
código, Consola, Environment para ver lo usado actualmente, archivos,
graficas, y otras más.

\textbf{¿Dónde irías para escribir instrucciones?} En Script para
guardarlas o en consola para ejecutarlas de inmediato.

\textbf{¿Dónde aparecen los resultados?} Texto y tablas en Consola;
gráficos en la pestaña Plots.

\textbf{¿Cómo explicarías esto a alguien que nunca ha programado?} Como
cocinar: receta (script), estufa (consola), despensa (environment) y un
celular para tomar las fotos o investigar (plots/help).

\subsection{2. Juega con R paso a paso:}\label{juega-con-r-paso-a-paso}

\subsection{3. Toma una captura de pantalla y pégala en el documento,
mostrando que ejecutaste el código en
RStudio.}\label{toma-una-captura-de-pantalla-y-puxe9gala-en-el-documento-mostrando-que-ejecutaste-el-cuxf3digo-en-rstudio.}

\begin{Shaded}
\begin{Highlighting}[]
\CommentTok{\# Cargar el dataset}
\FunctionTok{data}\NormalTok{(mtcars)}
\CommentTok{\# Ver las primeras filas}
\FunctionTok{head}\NormalTok{(mtcars)}
\end{Highlighting}
\end{Shaded}

\begin{verbatim}
##                    mpg cyl disp  hp drat    wt  qsec vs am gear carb
## Mazda RX4         21.0   6  160 110 3.90 2.620 16.46  0  1    4    4
## Mazda RX4 Wag     21.0   6  160 110 3.90 2.875 17.02  0  1    4    4
## Datsun 710        22.8   4  108  93 3.85 2.320 18.61  1  1    4    1
## Hornet 4 Drive    21.4   6  258 110 3.08 3.215 19.44  1  0    3    1
## Hornet Sportabout 18.7   8  360 175 3.15 3.440 17.02  0  0    3    2
## Valiant           18.1   6  225 105 2.76 3.460 20.22  1  0    3    1
\end{verbatim}

\begin{Shaded}
\begin{Highlighting}[]
\CommentTok{\# Ver los nombres de las columnas}
\FunctionTok{colnames}\NormalTok{(mtcars)}
\end{Highlighting}
\end{Shaded}

\begin{verbatim}
##  [1] "mpg"  "cyl"  "disp" "hp"   "drat" "wt"   "qsec" "vs"   "am"   "gear"
## [11] "carb"
\end{verbatim}

\begin{Shaded}
\begin{Highlighting}[]
\CommentTok{\# Calcular la media de las millas por galón}
\FunctionTok{mean}\NormalTok{(mtcars}\SpecialCharTok{$}\NormalTok{mpg)}
\end{Highlighting}
\end{Shaded}

\begin{verbatim}
## [1] 20.09062
\end{verbatim}

\begin{Shaded}
\begin{Highlighting}[]
\CommentTok{\# Filtrar autos con más de 6 cilindros}
\FunctionTok{subset}\NormalTok{(mtcars, cyl }\SpecialCharTok{\textgreater{}} \DecValTok{6}\NormalTok{)}
\end{Highlighting}
\end{Shaded}

\begin{verbatim}
##                      mpg cyl  disp  hp drat    wt  qsec vs am gear carb
## Hornet Sportabout   18.7   8 360.0 175 3.15 3.440 17.02  0  0    3    2
## Duster 360          14.3   8 360.0 245 3.21 3.570 15.84  0  0    3    4
## Merc 450SE          16.4   8 275.8 180 3.07 4.070 17.40  0  0    3    3
## Merc 450SL          17.3   8 275.8 180 3.07 3.730 17.60  0  0    3    3
## Merc 450SLC         15.2   8 275.8 180 3.07 3.780 18.00  0  0    3    3
## Cadillac Fleetwood  10.4   8 472.0 205 2.93 5.250 17.98  0  0    3    4
## Lincoln Continental 10.4   8 460.0 215 3.00 5.424 17.82  0  0    3    4
## Chrysler Imperial   14.7   8 440.0 230 3.23 5.345 17.42  0  0    3    4
## Dodge Challenger    15.5   8 318.0 150 2.76 3.520 16.87  0  0    3    2
## AMC Javelin         15.2   8 304.0 150 3.15 3.435 17.30  0  0    3    2
## Camaro Z28          13.3   8 350.0 245 3.73 3.840 15.41  0  0    3    4
## Pontiac Firebird    19.2   8 400.0 175 3.08 3.845 17.05  0  0    3    2
## Ford Pantera L      15.8   8 351.0 264 4.22 3.170 14.50  0  1    5    4
## Maserati Bora       15.0   8 301.0 335 3.54 3.570 14.60  0  1    5    8
\end{verbatim}

\subsection{4.Analizando}\label{analizando}

\textbf{¿Qué crees que fue más fácil?} Cargar mtcars y verlo con head(),
listar columnas con colnames(), y calcular la media con
mean(mtcars\$mpg).

\textbf{¿Qué resultó confuso?} Saber que mtcars ya está precargado en R,
entender \$ para acceder a columnas, y usar subset() con lógica
booleana.

\textbf{¿Qué errores te imaginas que cometerían tus estudiantes?}
Escribir mal nombres, olvidar mayúsculas, confundir = con ==, omitir el
dataset en las funciones.

\textbf{¿Cómo enseñarías estas primeras instrucciones?} Comparar con
Excel (filas=autos, columnas=variables, \$=columna llamada\ldots),
mostrar head() y View(), avanzar paso a paso, quizá dar un ejemplo y
luego dejarlos practicar.

\textbf{¿Qué analogías, metáforas o ejemplos usarías?} \$=etiqueta de
cajón, subset()=filtro de Excel, dataset=libreta de contactos.

\textbf{¿Qué herramientas visuales o estrategias crees que pueden
ayudar?} Usar View(), diagramas de tabla, comparaciones con Excel,
ejercicios cortos y colores para resaltar elementos clave.

\subsection{5. Profundicemos en la construcción de estas habilidades
docentes.}\label{profundicemos-en-la-construcciuxf3n-de-estas-habilidades-docentes.}

Abre RStudio y crea un nuevo script:

\subsection{6. Escribe y ejecuta el siguiente código en tu
script:}\label{escribe-y-ejecuta-el-siguiente-cuxf3digo-en-tu-script}

\begin{Shaded}
\begin{Highlighting}[]
\CommentTok{\# Filtrar autos con más de 6 cilindros}
\NormalTok{autos\_mas6c }\OtherTok{\textless{}{-}} \FunctionTok{subset}\NormalTok{(mtcars, cyl }\SpecialCharTok{\textgreater{}} \DecValTok{6}\NormalTok{)}
\CommentTok{\# Guardar el nuevo dataframe como CSV}
\FunctionTok{write.csv}\NormalTok{(autos\_mas6c, }\StringTok{"autos\_mas6c.csv"}\NormalTok{)}
\end{Highlighting}
\end{Shaded}

\subsection{7. Toma una captura de pantalla donde se ve el código y los
resultados en la consola de R y cópiala en el documento. (paso
3)}\label{toma-una-captura-de-pantalla-donde-se-ve-el-cuxf3digo-y-los-resultados-en-la-consola-de-r-y-cuxf3piala-en-el-documento.-paso-3}

\subsection{8. Observa y reflexiona, mientras ejecutas el código,
después responde en el
documento}\label{observa-y-reflexiona-mientras-ejecutas-el-cuxf3digo-despuuxe9s-responde-en-el-documento}

\textbf{¿Qué parte te pareció más clara y cuál podría confundir a un
estudiante sin experiencia?} La parte más clara es subset() si se
compara con un filtro en Excel, y lo más confuso sería entender que
\textless- crea un nuevo objeto y que write.csv() guarda un archivo en
la computadora.

\textbf{¿Qué conceptos requieren una explicación previa?} Sería
necesario explicar qué es un data frame, cómo funciona subset() para
filtrar filas, el operador \textless- para asignar resultados y
write.csv() para exportar datos.

\textbf{¿Qué metáforas o comparaciones podrías usar para explicar esos
conceptos?} \textless- sería como poner el resultado en un nuevo plato,
subset() como usar un colador para quedarte solo con lo que cumpla una
condición, un data frame como una hoja de Excel, y write.csv() como
guardar esa hoja para compartirla.

\subsection{9. Conecta con tu rol docente y completa el documento
realizando lo
siguiente:}\label{conecta-con-tu-rol-docente-y-completa-el-documento-realizando-lo-siguiente}

data(mtcars) Le pedimos a R que cargue un conjunto de datos que ya viene
instalado, como abrir un ejemplo que viene con el programa. head()
Muestra solo las primeras filas de la tabla para revisarla sin
abrumarnos; aquí pedimos 5 para reconocer su contenido.
mean(mtcars\$mpg) Calcula el promedio de la columna mpg; el signo de
pesos indica que queremos esa columna específica. subset(mtcars, cyl
\textgreater{} 6) Filtra las filas que cumplen la condición de tener más
de 6 cilindros. Guarda el resultado filtrado en un nuevo objeto llamado
autos\_mas6c. write.csv() Exporta el nuevo conjunto de datos a un
archivo CSV que se puede abrir o compartir.

Ideas para introducir el concepto de data frame - Hacer una tabla de los
estudiantes con diferentes datos, como año de nacimiento, película
favorita, que elijan entre 3 opciones una, y después mencionar que se
puede hacer con cada tipo de dato.

Sugerencias de errores frecuentes y cómo acompañar al estudiante Hacer
errores al escribir código para mostrar cómo solucionarlos, ver que
ocurre puede grabarles el ejemplo: -Escribir mal el nombre del objeto
(mtcars como mtcar): mostrar cómo usar autocomplete o revisar
Environment. -Olvidar mayúsculas o minúsculas: reforzar que R es
case-sensitive. -Confundir = con == al comparar: repasar diferencias
entre asignar y comparar. -No encontrar el archivo CSV: guiar para
revisar y cambiar el directorio de trabajo con getwd() / setwd(). Al
final entregar un archivo con errores frecuentes que ocurren.

Propuesta visual o actividad breve Usar la tabla de ejemplo y pedir a
los estudiantes que señalen una columna, calculen un promedio y filtren
por una condición. Usar los datos para hacer unas gráficas simples y que
vean que es más sencillo interpretarlos.

\subsection{10. Finalmente, sube al LMS:}\label{finalmente-sube-al-lms}

-Hice un .Rmd que contendría el código y las anotaciones, más parecido a
lo que aprendimos en el curso -Se me hace raro hacer el R, porque pues
solo era copiar

\end{document}
